\chapter{Амальгамированное объединение KO6 и семейно-калибровочного следа}

Модулярная конструкция на \(M_{18}(\mathbb C)\) решила задачу ориентации,
а старый носитель \(M_{60}(\mathbb C)\) решил задачу единого следа для
семейного и калибровочного чтений. Их обычное тензорное произведение имело
бы размерность \(18\cdot60=1080\) и дважды учитывало бы поколенческий
триплет. Поэтому проверяется склейка по общей физической части, а не
произведение независимых моделей.

\section{Общий 45-мерный угол}

Частичный угол старого меню имеет вид
\begin{equation}
 \mathcal H_{60}^{(p)}
 =\mathbb C^4_{\mathrm{menu}}\otimes
 (\mathbf{10}\oplus\overline{\mathbf5}),
 \qquad \dim\mathcal H_{60}^{(p)}=60.
\end{equation}
После канонического разложения \(\mathbb C^4=\mathbf1\oplus\mathbf3\)
физический семейный угол равен
\begin{equation}
 \mathcal H_{45}
 =\mathbb C^3_{\mathrm{fam}}\otimes
 (\mathbf{10}\oplus\overline{\mathbf5}),
 \qquad \dim\mathcal H_{45}=45.
\end{equation}

Распространение пакета \(\mathbf{10}\oplus\overline{\mathbf5}\) вдоль
трёх вершин KO6 даёт частичную цепь
\begin{equation}
 \mathcal H_{135}^{(p)}
 =(V_L\oplus V_G\oplus V_R)
 \otimes(\mathbf{10}\oplus\overline{\mathbf5}),
 \qquad \dim\mathcal H_{135}^{(p)}=135.
\end{equation}
Её левый угол \(V_L\otimes(\mathbf{10}\oplus\overline{\mathbf5})\)
отождествляется именно с \(\mathcal H_{45}\). Поэтому частичная размерность
склейки равна
\begin{equation}
 60+135-45=150,
\end{equation}
а после вещественного удвоения
\begin{equation}
 \dim\mathcal H_{\mathrm{par}}=300.
 \label{eq:v5-m300-dimension}
\end{equation}
Таким образом, \(780\) состояний наивного тензорного произведения не
появляются.

\section{Одна полная матричная алгебра}

Родительским операторным носителем служит
\begin{equation}
 \mathcal B_{\mathrm{par}}=M_{300}(\mathbb C),
 \qquad
 \tau_{300}(A)=\frac1{300}\Tr_{300}A.
\end{equation}
Он прост и имеет единственный нормированный след. Удвоенный угол меню имеет
ранг (120), угол цепи --- ранг (270), а их пересечение --- ранг (90):
\begin{equation}
 120+270-90=300.
\end{equation}
Следовательно, свободного центрального веса между двумя исходными
конструкциями нет.

На частичном углу меню физический проектор имеет ранг (45) из (60),
то есть
\begin{equation}
 \frac{45}{60}=\frac34.
\end{equation}
Внутри него тяжёлый семейный проектор по-прежнему имеет ранг (15), поэтому
\begin{equation}
 \frac{15}{45}=\frac13.
\end{equation}
Оба старых условных следа \(M_{60}\) и \(M_{45}\) тем самым являются
ограничениями одного следа \(M_{300}\), а не независимо выбранными мерами.

\section{Почему пакет SU(5) должен идти по всей цепи}

Если пакет \(\mathbf{10}\oplus\overline{\mathbf5}\) поместить только на
семейной вершине (L), ненулевое ребро не будет коммутировать с действием
SU(5). Для явного диагонального генератора машинная норма коммутатора равна
\(2.8752\). Если один и тот же пакет действует на всех трёх вершинах,
\begin{equation}
 D_{135}=D_9\otimes I_{15},
 \qquad
 T_{SU(5)}=I_9\otimes T_{15},
\end{equation}
то
\begin{equation}
 [D_{135},T_{SU(5)}]=0
\end{equation}
точно.

Это не вводит новую аномалию. Каждый пакет
\(\mathbf{10}\oplus\overline{\mathbf5}\) имеет нулевую кубическую аномалию
SU(5), а после разложения также нулевые смешанные и кубическую
гиперзарядовую аномалии. Максимальный машинный остаток равен
\(2.0\cdot10^{-16}\).

\section{Лёгкое ядро и массивные связующие комбинации}

На радиальном срезе положим
\begin{equation}
 X=\rho I_3,
 \qquad Y=rI_3.
\end{equation}
Тогда частичный оператор цепи имеет спектр
\begin{equation}
 \Spec D_{135}
 =\left\{
 0^{(45)},
 +\sqrt{\rho^2+r^2}^{\,(45)},
 -\sqrt{\rho^2+r^2}^{\,(45)}
 \right\}.
 \label{eq:v5-m300-chain-spectrum}
\end{equation}
Для проверочных значений \(\rho=0.83\), \(r=0.61\) ненулевая масса равна
\(1.03004854\). Следовательно, распространение пакета SU(5) по трём
вершинам не оставляет девять независимых лёгких поколений. Оно оставляет
одну 45-мерную линейную комбинацию, соответствующую трём поколениям, а две
ортогональные 45-мерные комбинации становятся массивными.

Для случайной невырожденной матрицы \(X\) и ненулевого скалярного \(Y\)
внутреннее ядро имеет размерность три. После тензорного умножения на пакет
SU(5) это снова даёт ровно (45) лёгких состояний.

\begin{theorem}[Отсутствие двойного счёта лёгких поколений]
Амальгамированная цепь имеет ровно 45-мерное частичное ядро при
невырожденном \(X\) и ненулевом \(Y\). Два дополнительных вершинных пакета
образуют массивные комбинации, а не дополнительные лёгкие поколения.
\end{theorem}

\section{Общая нормировка двух секторов}

На полном среднем секторе ранг равен
\begin{equation}
 2\cdot3\cdot15=90.
\end{equation}
Поэтому один след \(M_{300}\) даёт
\begin{equation}
 \frac1{90}\Tr\left(
 \widehat p_G[d,d^\dagger]^2\otimes I_{15}
 \right)
 =\tau_3(XX^\dagger-Y^\dagger Y)^2.
\end{equation}
Дефект равен \(1.8\cdot10^{-15}\). В том же родительском носителе
трёхпоколенческие индексы подгрупп SU(5) остаются
\begin{equation}
 I_{SU(3)}=I_{SU(2)}=I_{U(1)}=6.
\end{equation}
Таким образом, семейное отображение момента и калибровочные индексы
нормируются одним следом без относительного коэффициента.

\section{Граница кинематического успеха}

Получен первый кандидат на общую кинематическую родительскую архитектуру:
\begin{equation}
 \boxed{
 \mathcal B_{\mathrm{par}}=M_{300}(\mathbb C),
 \quad
 \mathcal H_{\mathrm{par}}=\mathbb C^{300}
 }
\end{equation}
содержит как углы модулярной KO6-цепи, так и старые углы \(M_{60}\) и
\(M_{45}\), причём все их следы индуцируются единственным
\(\tau_{300}\).

Однако спектральное равенство~\eqref{eq:v5-m300-chain-spectrum} ещё не
доказывает устойчивость вакуума. Остаются четыре задачи:
\begin{enumerate}
 \item вывести одно действие, фиксирующее \(X,Y\) и их относительную
 жёсткость;
 \item показать положительность полного физического гессиана массивных
 связующих комбинаций;
 \item обосновать удаление 15-мерного равномерного опорного сектора;
 \item включить семейный проектор ранга один и юкавские операторы в то же
 действие без нового коэффициента.
\end{enumerate}

\begin{nogocriterion}[Критерий провала действия \(M_{300}\)]
Кандидат закрывается, если полный симметрийно допустимый функционал требует
независимого относительного коэффициента между калибровочным и семейным
секторами, оставляет дополнительную нулевую моду либо делает одну из двух
связующих комбинаций тахионной.
\end{nogocriterion}

\section{Вердикт}

\begin{protocol}
\begin{enumerate}
 \item склейка без двойного счёта:
 \textbf{выполнена};
 \item один простой родительский носитель \(M_{300}\):
 \textbf{получен};
 \item один след для семейного и калибровочного секторов:
 \textbf{получен};
 \item представление SU(5) и компенсация аномалий:
 \textbf{сохранены};
 \item ровно 45 лёгких частичных состояний:
 \textbf{получены};
 \item точная модулярная ориентация и отображение момента:
 \textbf{сохранены};
 \item кинематический кандидат родительской архитектуры:
 \textbf{получен};
 \item полное действие и физический гессиан:
 \textbf{не получены};
 \item физическое замыкание:
 \textbf{не достигнуто}.
\end{enumerate}
\end{protocol}

Следующая проверка должна построить наиболее общий симметрийно допустимый
функционал \(M_{300}\) и вычислить его полный гессиан до феноменологии.

Машинная проверка и сертификат сохранены в
\path{s2t_v5_modular_ko6_m60_amalgamation_gate.py} и
\path{s2t_v5_modular_ko6_m60_amalgamation_gate_results.json}.